\documentclass[11pt]{scrartcl}
\usepackage[utf8]{inputenc}
\usepackage[sexy]{evanchen} 
\usepackage{graphicx}
\graphicspath{{8 Outsource K Math/Gambar/}}


\title{Barisan dan Deret}
\author{Azzam and Akeyla}
\date{Januari 2025}

\begin{document}

\maketitle
\section{Barisan Aritmetika, Geometri, dan Aritmetika-Geometri}
\subsection{Barisan Aritmetika}
Barisan aritmetika adalah barisan yang setiap suku yang bersebelahannya memiliki selisih yang tetap. Secara umum, barisan aritmetika berbentuk $a, a + b, a + 2b, ...$.\\
Suku ke-$n$, $U_n$ dapat dirumuskan dengan $U_n  = a + (n-1)b$.\\
Jumlah $n$ suku pertama, $S_n$, dapat dirumuskan dengan
\[S_n  = \dfrac{n}{2} (a + U_n )\]

\subsection{Barisan Geometri}
Barisan geometri adalah barisan yang setiap suku yang bersebelahannya memiliki rasio yang tetap. Secara umum, barisan geometri berbentuk $a,ar,ar^2, ...$.\\
Suku ke-$n$, $U_n$ dapat dirumuskan dengan $U_n = ar^{n – 1}$.\\
Jumlah $n$ suku pertama, $S_n$, dapat dirumuskan dengan \[S_n  = a\left(\dfrac{r^n - 1}{r - 1}\right).\]
Jika $-1<r<1$, maka jumlah barisan geometri takhingga bisa diperoleh, yakni
\[S_\infty = \dfrac{a}{1 - r}.\]

\subsection{Barisan Aritmetika-Geometri}
Merupakan barisan yang setiap suku yang bersebelahannya merupakan perkalian suku dari barisan aritmetika dan geometri. Salah satu contoh barisan ini adalah \[\dfrac{1}{2}, \dfrac{2}{2^2}, \dfrac{3}{2^3}, \dfrac{4}{2^4}, ...\]
Ketika rasio dari barisan geometri bernilai di rentang $(-1, 1)$, penjumlahan suku barisan aritmetika-geometri tak hingga dapat dihitung. Sebagai contoh,
\begin{align*}
    S &= \dfrac{1}{2} + \dfrac{2}{2^2} + \dfrac{3}{2^3} + \dfrac{4}{2^4} + ... \\
    \dfrac{1}{2}S &= \hspace{20pt} \dfrac{1}{2^2} + \dfrac{2}{2^3} + \dfrac{3}{2^4} + ... \\
\end{align*}
Dengan mengurangkan kedua nilai tersebut, diperoleh nilai $\dfrac{1}{2}S$ adalah sebesar \[\dfrac{1}{2} + \dfrac{1}{2^2} + \dfrac{1}{2^3} + \dfrac{1}{2^4} + ...\]
Deret yang terbentuk merupakan deret geometri yang bernilai \[\dfrac{1}{2}S = \dfrac{\dfrac{1}{2}}{1 - \dfrac{1}{2}} = 1\]
Dan dapat disimpulkan bahwa $S = 2$.

\section{Barisan Rekursif}
Suatu barisan dikatakan rekursif apabila nilai suku di barisan bergantung dari nilai suku-suku sebelumnya. Sebagai contoh, barisan Fibonacci yang didefinisikan sebagai $F_0 = 0, F_1 = 1,$ dan \[F_{n + 2} = F_{n + 1} + F_n\] untuk setiap bilangan cacah $n$.\\
Contoh lainnya adalah barisan Catalan yang didefinisikan dengan $C_0 = 1$ dan \[C_{n+1} = \sum_{i = 0}^n C_iC_{n-i}\]
Beberapa barisan dapat dinyatakan dalam bentuk eksplisit, yakni nilai suku ke-$n$ hanya berdasarkan dari nilai $n$. Sebagai contoh, nilai suku Catalan ke-$n$ adalah \[\dfrac{1}{n + 1} \binom{2n}{n}\]
Namun, tidak semua barisan rekursif dapat dinyatakan dalam bentuk eksplisit yang bagus.

\section{Periodisitas}
Suatu barisan $\{a_n\}$ dikatakan periodik dengan periode $k$ apabila $a_{n + k} = a_n$ untuk setiap bilangan asli $n$ yang cukup besar.\\
Apabila kita dapat menemukan bahwa sebuah barisan bersifat periodik, kita dapat mengetahui nilai suku dengan indeks yang besar dari nilai suku yang kecil.

\section{Teleskoping}

Prinsip teleskoping banyak digunakan dalam menentukan jumlahan atau perkalian suatu deret. Bentuk yang umum digunakan adalah:

\textbf{Teleskoping Penjumlahan:} \[\sum_{i=1}^n \left(a_{i+1} - a_i\right) = (a_2 - a_1) + (a_3 - a_2) + ... + (a_{n+1} - a_n) = a_{n+1} - a_1 \]

\textbf{Teleskoping Perkalian:} \[\prod_{i=1}^n \left(\frac{a_{i+1}}{a_i}\right) = \frac{a_2}{a_1} \cdot \frac{a_3}{a_2} \cdot ... \cdot \frac{a_{n+1}}{a_n} = \frac{a_{n+1}}{a_1}\]

Mari kita bahas salah satu contoh klasik:
\begin{example}
Tentukanlah nilai dari \[\frac{1}{1\cdot 2} + \frac{1}{2\cdot 3} + ... + \frac{1}{n(n+1)}\]
\end{example}

\begin{proof}[Solusi]
Perhatikanlah bahwa 
\begin{align*}
    \frac{1}{1\cdot 2} + \frac{1}{2\cdot 3} + ... + \frac{1}{n(n+1)} &= \sum_{i=1}^{n} \left(\frac{1}{k(k+1)}\right)\\
    &= \sum_{i=1}^{n} \left(\frac{1}{k} - \frac{1}{k+1}\right)\\
    &= \frac{1}{1} - \frac{1}{n+1}\\
    &= \frac{n}{n+1}
\end{align*}
\end{proof}

\begin{example}
    Tentukan nilai dari $1!1 + 2!2 + 3!3 + ... + n!n$
\end{example}

\begin{proof}[Solusi]
    Perhatikan bahwa 
    \begin{align*}
        k!k &= k!(k + 1 - 1)\\
        &= k!(k + 1) - k!\\
        &= (k + 1)! - k!
    \end{align*}
    Sehingga, 
    \begin{align*}
        &1!1 + 2!2 + 3!3 + ... + n!n \\
        = &(2! - 1!) + (3! - 2!) + ... + ((n + 1)! - n!) \\
        = &(n + 1)! - 1
    \end{align*}
\end{proof}


\section{Tips dan Trik}
\begin{enumerate}
    \item Coba cek untuk suku dengan indeks kecil terlebih dahulu
    \item Jika menemukan pola pada suku dengan indeks kecil, coba lakukan induksi agar dugaan kita terbukti
    \item Cek juga apakah barisan tersebut periodik
\end{enumerate}


\section{Latihan}
\begin{problem}
    (OSK 2021 Kemampuan Dasar) Misalkan $u_1, u_2, u_3, ...$ barisan aritmetika dengan suku-suku real positif. Jika $\dfrac{u_1 + u_2}{u_3} = \dfrac{11}{21}$, maka nilai dari $\dfrac{u_2 + u_3}{u_1}$ adalah ...
\end{problem}

\begin{problem}
    (AIME 1984) Tentukan nilai dari $a_2 + a_4 + a_6 + \dots + a_{98}$ apabila $a_1$, $a_2$, $a_3$, $\dots$ merupakan sebuah barisan aritmetika dengan selisih $1$ dan $a_1 + a_2 + a_3 + \dots + a_{98} = 137$.
\end{problem}

\begin{problem}
    (Turki MO, Penyisihan Pertama 2006) Barisan $(a_n)$ didefinisikan sebagai $a_1 = -1, a_2 = 2,$ dan \[a_n = \frac{a_{n-1}}{a_{n-2}}\] untuk $n\geq 3$. Nilai dari $a_{2021}$ adalah ....
\end{problem}

\begin{problem}
    (AIME 1988) Untuk setiap bilangan asli $k$, notasikan $f_1(k)$ sebagai kuadrat dari jumlah digit di $k$. Untuk $n \geq 2$, misalkan $f_n(k) = f_1(f_{n - 1}(k))$. Tentukan nilai $f_{1988}(11)$.
\end{problem}

\begin{problem}
    Tentukan nilai dari \[\dfrac{1}{1\cdot 3} + \dfrac{1}{3 \cdot 5} + ... + \dfrac{1}{99\cdot 101}\]
\end{problem}

\begin{problem}
    Nilai dari \[\left(1 - \dfrac{1}{2^2}\right)\left(1 - \dfrac{1}{3^2}\right)...\left(1 - \dfrac{1}{100^2}\right)\] adalah ...
\end{problem}

\begin{problem}
    (OSK 2022 Kemampuan Dasar) Untuk setiap bilangan asli $n$, misalkan $S(n)$ adalah jumlah dari semua digit-digit dari $n$. Diberikan barisan $\{a_n\}$ di mana $a_1 = 5$ dan $a_n = \left(S\left(a_{n-1}\right)\right)^2 - 1$ untuk $n \geq 2$. Sisa pembagian $a_1 + a_2 + ... + a_{2022}$ dengan $21$ adalah ...
\end{problem}

\begin{problem}
    (OSK 2022 Kemampuan Lanjut) Jika \[\sum_{k = 1}^{\infty} \dfrac{2k + B}{3^{k + 1}} = 10, \]
    maka nilai dari $B$ adalah ...
\end{problem}

\begin{problem}
    Tentukan nilai dari \[\dfrac{1}{2^1} + \dfrac{1}{2^2} + \dfrac{2}{2^3} + \dfrac{3}{2^4} + \dfrac{5}{2^5} + \dfrac{8}{2^6} + \dfrac{13}{2^7} + ...\]
\end{problem}

\begin{problem}
    Barisan $(x_n)$ memenuhi $x_0=3,x_1=4$, dan $x_{n+1}=x_{n-1}^2-nx_n$ untuk sembarang bilangan asli $n$. Tentukan nilai dari $\dfrac{x_0+x_1+x_2+...+x_{2022}}{2023}$
\end{problem}

\begin{problem}
    (Tournament of the Towns, Autumn 1997, Junior; OSP 2010-an) Barisan $(x_n)$ didefinisikan sebagai $x_1 = 20, x_2 = 21,$ dan \[x_{n+2} = x_n - \frac{1}{x_{n+1}}\] untuk $n\geq 1$. Tunjukkan bahwa terdapat bilangan asli $k$ sedemikian hingga $x_k = 0$ dan tentukan nilai $k$ tersebut. 
\end{problem}

\end{document}
